%%%%%%%%%%%%
%
% $Autor: Wings $
% $Datum: 2019-03-05 08:03:15Z $
% $Pfad: Specifications.tex $
% $Version: 4250 $
% !TeX spellcheck = en_GB/de_DE
% !TeX encoding = utf8
% !TeX root = ../../CPIManual
%
%%%%%%%%%%%%

\chapter{Specifications}

This chapter records the supported operating assumptions of the current dashboard release. It is intended as a compact technical reference for users who want to understand the scope of the application without reading backend source code.

\section{System Scope}

The application supports one forecasting use case: monthly CPI forecasting for Germany using the Destatis table \texttt{61111-0002}. The data is used as a univariate time series, and the implemented model inventory consists of ARIMA \texttt{(1,1,1)}, ETS \texttt{(A,A,A)}, and SARIMA \texttt{(1,1,1)$\times$(1,1,1,12)}.

\section{Input Specification}

\begin{itemize}
    \item \textbf{Primary input file:} \FILELOC{Code/data/61111-0002\_en.csv}
    \item \textbf{Expected frequency:} monthly
    \item \textbf{Expected target variable:} CPI index value with base year 2020=100
    \item \textbf{Observed usable date range in the current file:} January 1991 to February 2026
    \item \textbf{Usable observations after dropping future placeholders:} 422 monthly values
\end{itemize}

The data loader removes future placeholder entries such as \texttt{...}, maps month names to numeric values, and enforces a strict month-start frequency.

\section{Output Specification}

For normal dashboard use, the following prepared artifacts are expected to be present.

\begin{itemize}
    \item Serialized model files in \FILELOC{Code/saved\_models/}
    \item Plot images in \FILELOC{Code/plots/}
    \item Evaluation metrics in \FILELOC{Code/plots/evaluation\_metrics.csv}
    \item A LaTeX version of the metrics table in \FILELOC{Code/plots/evaluation\_metrics.tex}
\end{itemize}

\section{Supported User Controls}

The dashboard currently supports the following interactive controls and user actions.

\begin{itemize}
    \item single-model selection
    \item all-model comparison mode
    \item confidence interval level selection from 80\% to 99\%
    \item legend-based hiding and showing of plotted traces
    \item plot download through the built-in chart toolbar
\end{itemize}

The forecast horizon is not user-configurable in the present implementation. It is fixed to the 24-month test window used in both the training pipeline and the dashboard logic.

\begin{tcolorbox}[colback=blue!5!white, colframe=blue!75!black, title=Scope Reminder]
This dashboard is intentionally focused on one dataset and three classical models. It does \textbf{not} support arbitrary datasets, custom forecast horizons, or real-time data feeds. If you need those capabilities, refer to the Add Ons chapter for planned extensions.
\end{tcolorbox}

\section{Dependency Inventory}

The runtime dependencies declared in \FILELOC{Code/requirements.txt} are listed in Table~\ref{tab:manual-spec-deps}.

\begin{table}[H]
    \centering
    \small
    \setlength{\tabcolsep}{5pt}
    \renewcommand{\arraystretch}{1.15}
    \caption{Python dependencies required for operation.}
    \label{tab:manual-spec-deps}
    \begin{tabularx}{\textwidth}{|L{2.8cm}|X|}
        \hline
        \textbf{Package} & \textbf{Role in the project} \\
        \hline
        pandas & data loading, date parsing, indexing, and table handling \\
        \hline
        matplotlib & generation of saved forecast and decomposition plots \\
        \hline
        statsmodels & ARIMA, ETS, and SARIMA implementations \\
        \hline
        scikit-learn & RMSE, MAE, and MAPE evaluation metrics \\
        \hline
        jinja2 & dependency required by the current software environment \\
        \hline
        joblib & serialization and loading of trained model objects \\
        \hline
        streamlit & web application framework for the dashboard \\
        \hline
        plotly & interactive chart rendering in the dashboard \\
        \hline
    \end{tabularx}
\end{table}

\section{Artifact Flow Specification}

\begin{figure}[H]
    \centering
    \resizebox{0.95\textwidth}{!}{%
    \begin{tikzpicture}[node distance=1.2cm, >=latex, every node/.style={font=\small}]
        \tikzstyle{box}=[rectangle, rounded corners, draw=teal!60!black, fill=teal!10, align=center, text width=7.2cm, minimum height=0.9cm]
        \node[box] (csv) {1. Raw Destatis CSV stored in \FILELOC{Code/data/}};
        \node[box, below=of csv] (loader) {2. Data loader parses months, cleans values, and enforces monthly indexing};
        \node[box, below=of loader] (train) {3. Modeling pipeline trains ARIMA, ETS, and SARIMA and evaluates them on the 24-month holdout};
        \node[box, below left=1.25cm and 0.7cm of train] (artifacts) {4a. Save model artifacts, metrics, and plots};
        \node[box, below right=1.25cm and 0.7cm of train] (app) {4b. Streamlit app loads the data and the saved artifacts for presentation};

        \draw[->, thick] (csv) -- (loader);
        \draw[->, thick] (loader) -- (train);
        \draw[->, thick] (train) -- (artifacts);
        \draw[->, thick] (train) -- (app);
    \end{tikzpicture}%
    }
    \caption{Artifact flow from raw CPI data to the deployed dashboard.}
    \label{fig:spec-artifact-flow}
\end{figure}
